\subsection*{Restrictions}

\begin{itemize}

  
\item Re-declaration variables or functions is not allowed. Once a variable or function is defined, it cannot be redefined with the same name in the same scope
\footnote{
Scope refers to the region of a program in which a particular name (such as a variable, function, or class) is defined and can be accessed. In other words, 
it determines the part of the program where you can use that name without causing a name error. 
Scope is determined by the program's structure (usually its lexical or textual layout) and governs the visibility and lifetime of variables and other identifiers.

In \href{https://docs.python.org/3/tutorial/classes.html\#python-scopes-and-namespaces}{\color{DarkBlue}Python}, the \emph{scope} of a declaration is determined lexically: a variable declared inside a function is local to that function; 
if it is declared outside any function, it is global (i.e., module-level). Moreover, if a variable is declared in an enclosing function, 
it is available to inner functions (the enclosing scope), and if not found in these scopes, Python looks into the built-in scope.
}.

\item All names that are referenced in a program need to be declared, regardless of whether the
reference gets evaluated or not. A name is declared by an assignment, a function definition or a
parameter. The scope of a name that is declared in a function is that function, including the
functions nested in it; a name that is declared at the top level is in scope everywhere in the
program. A name that is declared in one function can therefore not be used in another function,
unless the second function is nested in the first. Built-in names are in scope everywhere. The
order of the declarations and the uses does not matter: a name that is used before it is
assigned is not rejected, but leads to a run-time error, as in Python.
\begin{lstlisting}
print(1 if True else x)
\end{lstlisting}
Leads to an error \texttt{name not found in the current or enclosing environment(s).}
\end{itemize}
